\documentclass[../../main.tex]{subfiles}
\begin{document}

\begin{flushright}
  \footnotesize\textit{【自動文字起こし・要確認】}
\end{flushright}

% Source: handwritten.pdf 第3問．
{\bf [解]}
題意のようになるのは，$n$回目に$n$のカードを引くか，$n$の1つ上におかれているカードを引く時．

\subparagraph{(i) $n$回目に$n$のカードを引く時}
各回の試行で，引いたカードをより下の位置におくので，
\begin{align}\dfrac{1}{n}\cdot\dfrac{2}{n}\cdot\dfrac{3}{n}\cdots\dfrac{n-1}{n}.
\end{align}
さらに，$n$のカードは1番上に戻るので，
\begin{align}
  \dfrac{1}{n}\cdot\dfrac{2}{n}\cdots\dfrac{n-1}{n}\cdot\dfrac{1}{n}
  =\dfrac{(n-1)!}{n^n}.
\end{align}

\subparagraph{(ii) $n$回目に$n$の1つ上におかれたカードを引く時}
$1$〜$n-1$回目の試行のうち1回だけ，$n$の上にカードをおく．これを$k$回目の試行として，
\begin{align}
  \dfrac{1}{n}\cdot\dfrac{2}{n}\cdots\dfrac{k-1}{n}\cdot\dfrac{n-k}{n}
  \cdot\dfrac{k}{n}\cdots\dfrac{n-2}{n}.
\end{align}
さらに，$n$回目に引いたカードは$n$のカードの下におくので，
\begin{align}
  \dfrac{1}{n}\cdot\dfrac{2}{n}\cdots\dfrac{k-1}{n}\cdot\dfrac{n-k}{n}
  \cdot\dfrac{k}{n}\cdots\dfrac{n-2}{n}\cdot\dfrac{n-1}{n}
  =\dfrac{(n-1)!(n-k)}{n^n}.
\end{align}
$k$について足して
\begin{align}
  \dfrac{(n-1)!}{n^n}\sum_{k=1}^{n-1}(n-k)
  =\dfrac{(n-1)!}{n^n}\sum_{k=1}^{n-1}k
  =\dfrac{(n-1)!}{n^n}\cdot\dfrac{1}{2}n(n-1).
\end{align}
\casesend
以上から
\begin{align}
  \dfrac{(n-1)!\left\{1+\dfrac{1}{2}n^2-\dfrac{1}{2}n\right\}}{n^n}
  =\dfrac{(n-1)!(n^2-n+2)}{2n^n}.
\end{align}

\end{document}
